\section{Results}
\label{sec:results}

\begin{table*}[t]
\caption{For lists with several correct choices, \emph{Gap} is the mean probability of the highest correct choice minus that of the lowest, in percentage points, \emph{Ranked} is the share of lists with every correct choice above every distractor, and \emph{Below $1/n$} is the share of those lists with a correct choice below 1/n for a list of n choices. $1/n = 0.1$ for every list except 39 HotpotQA lists of 3 to 9 passages. The two shares use the released order of the ECtHR lists and leave out the lists with no distractor.}
\label{tab:allocation}
\centering
\footnotesize
\begin{tabular}{lrrrrrrrrrrrr}
\toprule
 & \multicolumn{4}{c}{MetaTool} & \multicolumn{4}{c}{HotpotQA} & \multicolumn{4}{c}{ECtHR} \\
\cmidrule(lr){2-5} \cmidrule(lr){6-9} \cmidrule(lr){10-13}
Model & High & Low & Distr. & All $>$0.25 & High & Low & Distr. & All $>$0.25 & High & Low & Distr. & All $>$0.25 \\
\midrule
Jev & 0.85 & 0.11 & 0.04 & 15\% & 0.77 & 0.17 & 0.05 & 30\% & 0.72 & 0.15 & 0.07 & 23\% \\
Liquid d1 & 0.71 & 0.16 & 0.12 & 27\% & 0.74 & 0.19 & 0.05 & 35\% & 0.72 & 0.16 & 0.07 & 25\% \\
GPT-6 Luna Decisions & 0.66 & 0.11 & 0.17 & 16\% & 0.72 & 0.15 & 0.08 & 26\% & 0.73 & 0.13 & 0.06 & 17\% \\
\bottomrule
\end{tabular}

\end{table*}

\begin{table}[t]
\caption{Each list has 5 copies of one correct choice and is sent in its 5 cyclic permutations. The table shows the mean probability of the highest copy minus that of the lowest copy, in percentage points (\emph{Gap}). Equal probabilities give 0.}
\label{tab:copies}
\centering
\footnotesize
\setlength{\tabcolsep}{4pt}
\begin{tabular}{lrrrrrrrrr}
\toprule
 & \multicolumn{3}{c}{MetaTool} & \multicolumn{3}{c}{HotpotQA} & \multicolumn{3}{c}{ECtHR} \\
\cmidrule(lr){2-4} \cmidrule(lr){5-7} \cmidrule(lr){8-10}
Model & High & Low & First & High & Low & First & High & Low & First \\
\midrule
Jev & 0.32 & 0.11 & 36.2\% & 0.34 & 0.11 & 22.9\% & 0.29 & 0.12 & 3.6\% \\
Liquid d1 & 0.81 & 0.03 & 99.9\% & 0.73 & 0.04 & 99.5\% & 0.55 & 0.07 & 99.6\% \\
GPT-6 Luna Decisions & 0.78 & 0.02 & 98.6\% & 0.79 & 0.01 & 99.4\% & 0.67 & 0.02 & 98.9\% \\
\bottomrule
\end{tabular}

\end{table}

\begin{table}[t]
\caption{The share of answers that predict the copy in the first and in the last position, on the lists of Table~\ref{tab:copies}. An equal share is 20\%.}
\label{tab:positions}
\centering
\resizebox{\linewidth}{!}{\begin{tabular}{lrrrrrr}
\toprule
 & \multicolumn{2}{c}{MetaTool} & \multicolumn{2}{c}{HotpotQA} & \multicolumn{2}{c}{ECtHR} \\
\cmidrule(lr){2-3} \cmidrule(lr){4-5} \cmidrule(lr){6-7}
Model & First & Last & First & Last & First & Last \\
\midrule
Jev & 36\% & 25\% & 66\% & 16\% & 10\% & 56\% \\
Liquid d1 & 98\% & 1\% & 98\% & 1\% & 97\% & 1\% \\
GPT-6 Luna Decisions & 98\% & 1\% & 98\% & 1\% & 76\% & 1\% \\
\bottomrule
\end{tabular}
}
\end{table}

\begin{table}[t]
\caption{The highest exact set accuracy over the thresholds 0, 0.01, \ldots, 1. Exact set accuracy is the share of examples in which every correct choice is at or above the threshold and every distractor is below it. \emph{List} uses the probabilities of the list and \emph{Yes/no} the probability of yes of each choice.}
\label{tab:thresholds}
\centering
\resizebox{\linewidth}{!}{\begin{tabular}{lrrrrrr}
\toprule
 & \multicolumn{2}{c}{MetaTool} & \multicolumn{2}{c}{HotpotQA} & \multicolumn{2}{c}{ECtHR} \\
\cmidrule(lr){2-3} \cmidrule(lr){4-5} \cmidrule(lr){6-7}
Model & List & Yes/no & List & Yes/no & List & Yes/no \\
\midrule
Jev & 0.56 & 0.31 & 0.63 & 0.29 & 0.52 & 0.44 \\
Liquid d1 & 0.42 & 0.38 & 0.65 & 0.27 & 0.49 & 0.54 \\
GPT-6 Luna Decisions & 0.39 & 0.24 & 0.47 & 0.27 & 0.48 & 0.38 \\
\bottomrule
\end{tabular}
}
\end{table}

\paragraph{Correct choices receive unequal probabilities.}
When a list has two or more correct choices, the mean probability of the highest correct choice exceeds that of the lowest by 55 to 74 percentage points (Table~\ref{tab:allocation}, with 95\% intervals in Appendix~\ref{app:intervals}). A correct choice receives less than $1/n$ in 36\% to 60\% of the lists. The models rank every correct choice above every distractor in 52\% to 85\% of the lists. In 28\% to 50\% of those lists, a correct choice receives less than $1/n$. When the highest-probability correct choice is removed, the mean renormalized probability of the next correct choice differs from the mean probability the model assigns to that choice after the removal by 0.01 to 0.13 across models and datasets (Appendix~\ref{app:tables}). Correct choices are not equally relevant, and a probability below $1/n$ by itself is not a wrong allocation. However, relevance alone does not explain the size of the gap, as further demonstrated by the calibration and identical choices experiments below.

\paragraph{Identical copies receive unequal probabilities.}
In lists of five identical copies of a correct choice, the mean probability of the highest copy exceeds that of the lowest by 17 to 79 percentage points (Table~\ref{tab:copies}). Similar to \citet{deusser2026evaluating}, we observe that Jev and Liquid d1 show little position bias on lists with one correct tool, with accuracy changing by at most 4 percentage points between the first and the tenth place (Appendix~\ref{app:position}). However, identical copies of a correct choice still receive unequal probabilities. For Liquid d1 and GPT-6 Luna Decisions the gap between the highest and the lowest copy is 37 to 79 percentage points, and they predict the copy in the first position in 76\% to 98\% of their answers (Table~\ref{tab:positions}). For Jev it is 17 to 26 percentage points, and on ECtHR it predicts the copy in the last position in 56\% of its answers. Table~\ref{tab:copies-one} shows one HotpotQA list, in which Liquid d1 and GPT-6 Luna Decisions give 0.66 and 0.82 to the first copy. A choice can therefore receive a low probability for reasons unrelated to its content.

\paragraph{Thresholding.}
To give threshold-based selection the strongest possible advantage, we report the maximum exact set accuracy over all evaluated thresholds on each dataset. Even under this optimistic setting, exact set accuracy ranges from 0.34 to 0.64, despite the models ranking all correct choices above all distractors in 52\% to 85\% of the examples. For our lists of 10 choices each, the thresholds that maximize exact set accuracy range from 0.02 to 0.08 (Table~\ref{tab:thresholds}, Figure~\ref{fig:exact-set}). Thresholds this low have no practical interpretation, as they select choices to which the model gives a probability of a few percent. We repeat this process for the yes/no questions. The models rank all correct choices above all distractors in 42\% to 76\% of the examples, and the thresholds that maximize exact set accuracy range from 0.05 to 0.63. The highest, 0.63, is for Liquid d1 on ECtHR, where exact set accuracy is 0.54.

\paragraph{Calibration.}
When choice lists contain a single correct choice, the ECE is 1 to 14 percentage points. However, when lists with multiple correct choices are presented as yes/no questions, the probability of yes departs from the observed frequency in intermediate bins. While the ECE over all bins is 3 to 9 percentage points, it is influenced by low probabilities. Between 46\% and 78\% of the answers fall in the bin from 0 to 0.1, and at least 92\% of them are about distractors. Once this bin is removed, the ECE rises to 8 to 18 percentage points (Table~\ref{tab:calibration}, Figure~\ref{fig:reliability}). On intermediate bins, the gap between the observed frequency and the mean probability of yes reaches 19 percentage points for Jev and 40 percentage points on ECtHR. The probability of yes is therefore not calibrated with the observed frequency of correct choices.

\paragraph{No correct choice.}
When no choice in a list is correct, Jev and Liquid d1 answer every list, with a mean confidence of 0.49 to 0.87 (Appendix~\ref{app:abstention}).

\FloatBarrier
